\section{Lecture 18 (November 13th)}
\begin{rmk}
Last time we learned Galois theory. Today, we learn the main theorem and examples of $\mathrm{Gal}(F/K)$. 
\end{rmk}
\vspace{2ex}
\begin{defi}
(15.2) An algebraic extension $F/K$ is Galois if it is separable and normal.
\end{defi}
\vspace{2ex}
\begin{rmk}
A finite extension $F/K$ is Galois if and only if $F$ is splitting field of a separable polynomial in $K[t]$ if and only if 
\[[F:K]=|\mathrm{Aut}(F/K)|=|\mathrm{Gal}(F/K)|\]
\end{rmk}
\vspace{2ex}
\begin{rmk}
Let $F/K$ is finite
\[|\mathrm{Aut}(F/K)|\leq [F:K]_{s}\leq [F:K]\]
the first equality holds if and only if $F/K$ is normal and the second equality holds if and only if $F/K$ is separable. We see that the two conditions we previously sought were exactly $F$ being normal and separable. 
\end{rmk}
\vspace{2ex}
\begin{thm}
(15.31) and (15.35) (Fundamental theorem of Galois theory) Let $F/K$ be a finite Galois extension. Then the Galois correspondence
\[\{E:\mathrm{\ an\ intermediate\ field\ }K\subset E\subset F\}\longleftrightarrow \{\mathrm{ subgroups\ of\ }\mathrm{Gal}(F/K)\}\]
where $E\mapsto \mathrm{Gal}(F/E)$ and $H\mapsto F^{H}$ is a bijection.
\end{thm}
\vspace{2ex}
\begin{proof}
If $E$ is an intermediate field $K\subset E\subset F$, then the extension $E/K$ is Galois (or normal) if and only if $\mathrm{Gal}(F/E)$ is a normal subgroup of $\mathrm{Gal}(F/K)$. In this case, $\mathrm{Gal}(E/K)$ is isomorphic to $\mathrm{Gal}(F/K)/\mathrm{Gal}(F/E)$.
\end{proof}
\vspace{2ex}
\begin{ex}
(15.12) Let $f(t)=t^{4}+t^{3}+t^2+t+1\in {\bm Q}[t]$. Let $F=$ the splitting field of $f(t)$, equal to ${\bm Q}(\xi ,\xi ^2,\xi ^3,\xi ^{4})$. Here, $\xi $ is a primative 5th root of unity. Note that
\begin{itemize}
\item[(i)] $F/{\bm Q}$ is finite Galois.
\item[(ii)] Let $\sigma \in \mathrm{Gal}(F/{\bm Q})$. 
\begin{center}

\begin{tikzpicture}[>=stealth, node distance=3cm]

% Top nodes
\node (A) at (0,0) {${\bm Q}(\xi ,\xi ^2,\xi ^3,\xi^{4})$};
\node (B) at (6,0) {${\bm Q}(\xi ,\xi^2,\xi^3,\xi^{4})$};

% Arrow between them
\draw[->] (A) -- node[above] {$\sigma $} (B);

% Bottom node
\node (C) at (3,-3) {${\bm Q}$};

% Downward arrows
\draw[->] (A) -- (C);
\draw[->] (B) -- (C);

\end{tikzpicture}
\end{center}
\begin{itemize}
\item[(ii-i)] It is enough to consider $\sigma  (\xi )$ 
\item[(ii-ii)] $\sigma (\xi )$ must be zero of $f(t)$ ($f(\xi )=0$ implies that $f(\sigma (\xi ))=0$) which implies that $\sigma (\xi )=\xi ^{i}$ for $1\leq i \leq 4$. 
\item[(ii-iii)] There exists a $\sigma _{i}\in \mathrm{Gal}(F/K)$ such that $\sigma _{i}(\xi )=\xi ^{i}$ for each $1\leq i\leq 4$ by the isomorphism extension theorem. 
\item[(ii-iv)] $\mathrm{Gal}(F/{\bm Q})=\{\sigma_1,\sigma_2,\sigma_3,\sigma_4\}$
\item[(ii-v)] $({\bm Z}/5{\bm Z})^{\times }\rightarrow \mathrm{Gal}(F/{\bm Q})$.
\end{itemize}
\end{itemize}
\end{ex}
\vspace{2ex}

